\documentclass{beamer}
% Constructs that real decks use inside frames and that canvas editing must
% keep valid (see the corpus survey): written for the editing fuzzers.
\usepackage{amsmath,amssymb,booktabs}
\newcommand{\R}{\mathbb{R}}
\newcommand{\highlight}[1]{\textcolor{red}{#1}}
\newcommand{\remark}[2][Note]{\textbf{#1:} #2}
\DeclareMathOperator*{\argmin}{arg\,min}
\begin{document}

\begin{frame}{Line breaks}
Alpha beta gamma \\
delta epsilon\\[4pt]
zeta eta\newline theta iota \\
kappa.

Lambda mu nu.
\end{frame}

\begin{frame}{Item labels}
\begin{itemize}
  \item[] Unlabelled alpha.
  \item[$\to$] Arrow beta.
  \item[--] Dash gamma.
  \item[\textcolor{red}{A:}] Coloured delta.
\end{itemize}
\begin{description}
  \item[Term] Its description.
  \item[Other term] Another one.
\end{description}
\end{frame}

\begin{frame}{Item overlays}
\begin{itemize}[<+->]
  \item First alpha.
  \item
    Second, on its own line.
  \item<3-> Third with a specification.
  \begin{itemize}
    \item Nested beta.
    \item Nested gamma.
  \end{itemize}
  \pause
  \item Last delta.
\end{itemize}
\end{frame}

\begin{frame}{Overlay commands}
Alpha \uncover<2->{beta} \alert<3>{gamma} \alert<1>{delta}.

\textcolor<2->{blue}{Epsilon} \visible<2>{zeta} \alt<2>{eta}{theta}.

\onslide<2|handout:2>
Iota kappa.
\only<2->{Lambda mu.}
\pause
Nu xi.
\end{frame}

\begin{frame}{Only around items}
\begin{itemize}
  \item Alpha.
  \only<2->{
  \item Beta.
  \item Gamma.
  }
\end{itemize}
\end{frame}

\begin{frame}{Old-style groups}
Alpha {\bf beta} gamma {\em delta} and {\small epsilon
  zeta} eta.
\begin{itemize}
  \item {\bf Theta}: iota.
  \item \textcolor{red}{\bf Kappa} lambda.
\end{itemize}
$a + {\color{blue} b} = c$.
\end{frame}

\begin{frame}{Comments}
\begin{columns}%
\begin{column}{.5\textwidth}%
Alpha beta% glued
gamma. % trailing note
\end{column}%
\begin{column}{.5\textwidth}
\begin{itemize}
  \item Delta. % note
  % between items
  \item Epsilon.
  % \item Commented out.
  \item Zeta.
\end{itemize}
\end{column}
\end{columns}
\end{frame}

\begin{frame}[fragile]{Verbatim-like}
Use \verb|alpha_beta| or \texttt{gamma\_delta}.
See \url{https://example.org/~alpha#beta}.
\end{frame}

\begin{frame}{Rules and cells}
\begin{tabular}{|l|c|r|}
\hline
\multicolumn{2}{|c|}{Alpha beta} & 1\\
\hline
Gamma & & 2\\[2pt]
\cline{2-3}
{[0,1]} & Delta & 3\\
\hline
\end{tabular}
\end{frame}

\begin{frame}{Booktabs}
{\small
\begin{tabular}{lcc}
\toprule
Alpha & Beta & Gamma\\
\cmidrule(lr){2-3}
Delta & \textcolor{red}{\bf 1} & 2\\
\midrule
Epsilon & 3 & 4\\
\bottomrule
\end{tabular}}
\end{frame}

\begin{frame}{Math forms}
Inline $x^2 + y_i$ and \(f'(x) \le g''(x)\), sets $\{a \mid a > 0\}$ in $\R^n$.
\[
  \hat x = \argmin_{z \in \R^n} \|z\|, \quad
  \underbrace{x_1 + \dots + x_n}_{\text{terms}} \geq 0
\]
\begin{equation*}
  \mathbb{E}[X] = \sum_{k=1}^n k \, p_k
\end{equation*}
\end{frame}

\begin{frame}{Cases and alignments}
\[
  f(x) = \begin{cases} 0 & \text{if } x < 0,\\ x & \text{otherwise.} \end{cases}
\]
\begin{align*}
  \left. \frac{d}{dx} x^2 \right|_{x=1} &= 2 \notag\\
  \left\{ a, b \right\} &\subseteq S
\end{align*}
\end{frame}

\begin{frame}{Columns}
\begin{columns}[onlytextwidth,T]
\column{.4\textwidth}
Alpha beta.
\column{.6\textwidth}
\begin{itemize}
  \item Gamma.
\end{itemize}
\end{columns}
\begin{columns}[c]
\begin{column}{.5\textwidth}
Delta.
\end{column}
\begin{column}{.5\textwidth}
Epsilon.
\end{column}
\end{columns}
\end{frame}

\begin{frame}{Blocks}
\begin{block}{}
Untitled alpha.
\end{block}
\begin{block}<2->{Claim $1$}
Beta gamma.
\end{block}
\begin{alertblock}{\only<1>{Delta}\only<2>{Epsilon}}
Zeta.
\end{alertblock}
\begin{exampleblock}{Eta}
Theta.
\end{exampleblock}
\end{frame}

\begin{frame}{Links and notes}
An \emph{alpha} \href{https://example.org}{link}, see~\ref{beta} and \hyperlink{gamma}{delta}.
Epsilon\footnote{\url{https://example.org}, zeta.} eta.
\hypertarget<2>{gamma}{Theta.}
\href{https://example.org}{An \emph{alpha} link}
Alpha \textbf<1>{delta}.
\end{frame}

\begin{frame}{Spacing}
\vspace{-8pt}
Alpha\hfill{}beta.
\medskip

Gamma\hspace*{1em}delta\quad epsilon, e.g.\ zeta\ldots{} eta.
\bigskip
\vfill
\centerline{Theta}
\begin{center}
\small Iota.
\end{center}
\end{frame}

\begin{frame}{Typography}
``Alpha'' -- beta --- gamma: \& \% \$ \# \_ and~delta, vs.\ epsilon.
Accents \c{c} \'e \"o, direct é ü ł, and a citation-like [1].
\end{frame}

\begin{frame}[plain]
Plain alpha.
\end{frame}

\begin{frame}{}
Empty title beta.
\end{frame}

\begin{frame}<1->[label=frameheaders]{Title}{Subtitle}
Gamma.
\end{frame}

\begin{frame}
\frametitle{Title in the body}
Delta.
\pause
Epsilon.
\end{frame}

\begin{frame}[allowframebreaks]{Breaks}
Alpha beta.
\framebreak
Gamma delta.
\end{frame}

\begin{frame}{Boxes}
\mbox{alpha beta} \fbox{gamma} \colorbox{yellow}{delta} \makebox[3em][c]{epsilon}
\parbox{4em}{zeta\\ eta} \scalebox{.8}{theta} \rotatebox{10}{iota} \raisebox{2pt}{kappa}
\resizebox{2em}{!}{lambda}
\end{frame}

\begin{frame}{User macros}
Space $\R n$, \highlight{alpha}, \remark{beta} and \remark[Aside]{gamma}.
$\argmin_x f(x)$.
\end{frame}

\begin{frame}{Nested lists}
\begin{enumerate}
  \item Alpha.

  \item Beta.
  \begin{enumerate}
    \item Gamma.
    \begin{itemize}
      \item Delta.
      \item
    \end{itemize}
  \end{enumerate}
\end{enumerate}
\end{frame}

\begin{frame}
$x$ starts this body.
\end{frame}

\begin{frame}{Bracket first}
[Draft] alpha beta.
\end{frame}

\end{document}
